\scanpage{253}
\mspage{127}
% archive sheet number: 252

Soit donc $\deleted{\unsure{f|g}}\unsure{\varphi|f}:\deleted{X}I\to\deleted{Y}J$ dans $W$, avec $\pi_0(f)$ non injective. Procédant comme \unsure{tantôt|\deleted{tantôt}} en utilisant Loc $(\deleted{\unsure{6}}5)$, on est ramené au cas où de plus $\deleted{Y}J$ est $0$-connexe. (\uline{NB} dans les deux cas, on aurait pu aussi invoquer 9.6.) \deleted{On va utiliser les} Choisissons deux objets $\deleted{\unsure{x}}i_0,\deleted{\unsure{x}}i_1$ dans $\deleted{X}I$ dans deux composantes connexes distinctes, soient $\deleted{\unsure{y}}j_0,\deleted{\unsure{y}}j_1$ leurs images dans $\deleted{Y}J$. On va traiter $(I,i_0,i_1)$ et $(J,j_0,j_1)$ comme des ``intervalles homotopiques''.

Le fait que $i_0,i_1$ appartiennent à des composantes distinctes de $I$ implique que deux flèches $f_0,f_1:X\rightrightarrows Y$ sont ``homotopes par. à $(I,i_0,i_1)$'' i.e. il existe $h:X\times I\unsure{-|\to}Y$ dans $\mathrm{Cat}$ telle que $f_0=h\partial_0^X$, $f_1=h\partial_1^X$, où
\[
 \partial_0^X,\partial_1^X:X-I\times X
\]
sont donnés par $x\mapsto(\unsure{j|i}_\alpha\unsure{,|\times}x)$ $(\alpha\in\{0\,1\})$.

Pour trouver $h$, on décompose $\deleted{X}I$ en
\[
 \deleted{X}I=\deleted{X}I_0\amalg\deleted{X}I_1,\quad\text{avec } i_0\in I_0,\ i_1\in I,
\]
d'où
\[
 I\times \deleted{\unsure{I}}X\simeq(I_0\times \deleted{\unsure{I}}X)\amalg(I_1\times\unsure{X'|X)|X},
\]
et on définit $h$ par ses restrictions \inserted{$h_0,h_1$} aux deux morceaux $I_0\times X,I_1\unsure{<|\times}X$, savoir
\[
 h_0=\deleted{\unsure{\mathrm{pr}}}f_0\mathrm{pr}_2^{I_0},\quad
 h_1=f_1\circ\mathrm{pr}_2^{I_1},
\]
où
\[
 \mathrm{pr}_2^{I_\alpha}:I_\alpha\times X-X.
\]
